\label{sec:introduction}

\subsection{Section~203 in the VRA Framework}

The Voting Rights Act of 1965 is most widely discussed in redistricting
contexts through two provisions: Section~2, which prohibits voting
practices that dilute minority voting strength, and Section~5, which
required covered jurisdictions to seek federal preclearance before changing
voting procedures (now effectively dormant after
\textit{Shelby County v.\ Holder}\footnote{\textit{Shelby County
v.\ Holder}, 570 U.S.\ 529 (2013).}). But the VRA also contains
Section~203, a provision focused not on representation but on access:
requiring that voting materials be provided in the language of
qualifying minority groups so that language-minority citizens can
participate effectively in the electoral process.

Section~203, codified at 52~U.S.C.~\S~10503, was added to the VRA in 1975
and has been renewed several times, most recently in 2006 as part of the
Fannie Lou Hamer, Rosa Parks, and Coretta Scott King Voting Rights Act
Reauthorization and Amendments Act of 2006 (VRARA), Pub.\ L.\ No.\ 109-246,
120 Stat.\ 577. It requires any political subdivision where the conditions
of Section~203(b) are met to provide all election materials in both English
and the applicable language(s) of covered minority groups.

\subsection{Why Redistricting Affects Section~203 Administration}

Section~203 is determined at the jurisdictional level---counties and
equivalent units---based on Census Bureau determinations of language
minority population size and English literacy rates. When congressional
district boundaries cut through a Section~203-covered county, the
administrative requirements do not stop at the district line. The county
election office must provide language assistance throughout the county,
but the federal election administration (ballots, candidate information)
is organized at the congressional district level. This mismatch creates
coordination overhead: a county split between two congressional districts
must manage two separate sets of district-level materials in the required
language(s).

The coordination burden is greater when the split is not clean---when
the Section~203-covered language community is itself divided between
the two congressional districts. If the Vietnamese-speaking community
in Santa Clara County is split between two districts, each district's
representative has a reduced incentive to serve that community's
interests, and the community's political voice is diluted even though
the county as a whole retains Section~203 coverage.

\subsection{The \textsc{bisect} Connection}

\textsc{bisect} produces redistricting plans using a graph-partitioning
algorithm that treats census tracts as nodes and shared boundaries as
edges. The algorithm's county-weight mode (\texttt{--weights-override county})
assigns higher edge weights to the boundaries between tracts in different
counties, penalizing county splits in the optimization. This produces
plans with fewer county splits than algorithms that do not incorporate
county-boundary information---and that therefore tend to be more compatible
with Section~203 administration.

This paper quantifies the Section~203 administrative benefit of
\textsc{bisect}'s county-weight mode, compares it to enacted plans, and
analyzes the interaction between Section~203 coverage and Section~2
majority-minority district requirements.

\subsection{Overview}

Section~\ref{sec:background} sets out the Section~203 coverage formula,
the authorizing statute, and the 2023 coverage determination.
Section~\ref{sec:methodology} analyzes how redistricting affects
Section~203 administration, with a focus on county splits and their
administrative consequences. Section~\ref{sec:results} presents
empirical results for Texas and California.
Section~\ref{sec:legal} analyzes the interaction between Section~203
and Section~2 requirements. Section~\ref{sec:conclusion} concludes.
